% ajusta o espaçamento dos parágrafos do resumo
\setlength{\absparsep}{18pt} 


\begin{resumo}
    
Os modelos de linguagem de grande escala (LLMs) ampliaram as possibilidades de automação do atendimento ao cliente, permitindo interpretar solicitações, consultar informações e produzir respostas mais flexíveis do que as obtidas com fluxos baseados apenas em regras. Ainda não está claro, porém, em quais condições a distribuição das responsabilidades entre agentes especializados oferece vantagens suficientes  para compensar o aumento da complexidade e do consumo de recursos em relação a uma arquitetura de agente único. Este trabalho tem como objetivo comparar experimentalmente as duas arquiteturas em
um cenário controlado de atendimento textual de um provedor de Internet, considerando a correção e a relevância das respostas, a conclusão das tarefas, a precisão dos encaminhamentos, o tempo de resposta e o consumo de tokens. Para isso, serão implementados dois protótipos, submetidos ao mesmo conjunto
balanceado de casos de teste e às mesmas condições de execução, com avaliação apoiada pelo DeepEval e por métricas registradas diretamente durante o experimento. Espera-se identificar em quais situações a especialização por agentes setoriais melhora o atendimento e quais efeitos essa organização produz sobre a latência, o uso de tokens e o custo operacional.
    
\noindent\textbf{Palavras-chave:} modelos de linguagem de grande escala; sistemas multiagentes; chatbots; atendimento ao cliente; agentes inteligentes.

\end{resumo}



%-----------------------------------------------%
\begin{resumo}[Abstract]
\begin{otherlanguage*}{english}
    This is the english abstract.
\vspace{\onelineskip}

\noindent 
Keywords: latex; abntex; text editoration.
\end{otherlanguage*}
\end{resumo}
%-----------------------------------------------%